\begin{frame}
  \frametitle{Group Work}

  Individual grades, \(g\), are the product of your group's earned grade (\(r\), evaluated using the rubric) and your effort (\(f\), from ratings and commit history), that is, \(g = rf\).

  \pause
  \begin{block}{Example}
    \begin{enumerate}[<+->]
    \item You are part of a four person group which earned \(r = 90\%\).
    \item Each person is expected to contribute \(e = 1/4 = 25\%\).
    \item You and your teammates rated your contribution as \(w_{r} = 27.5\%\),\\ so your rated effort is \(f_{r} = w_{r} / e = 1.1\).
    \item Your total number of commits is similarly \(w_{c} = 27.5\%\),\\ thus your commit effort is \(f_{c} = w_{c} / e =  1.1\).
    \item Your effort, \(f = (w_{r} + w_{c}) / 2 = 1.1\).
    \item Thus your grade is \(g = rf = 99\%\).
    \end{enumerate}
  \end{block}

  % \note{
  %   \begin{itemize}
  %   \item You will always be given the values \(r\), \(e\), \(w_{r}\), and \(w_{c}\), and feedback from the rubric.
  %   \item Your self-reported rating is included in \(w_{r}\).
  %   \item \textbf{It is possible to \emph{loose} points from \(r\), by having contributed less.}
  %   \item I reserve the right to treat values of \(f\) near 1 as being exactly 1, to calculate \(w_{c}\) in several ways, and to weight \(w_{r}\) and \(w_{c}\) differently.
  %   \end{itemize}
  % }
  
\end{frame}

%%% Local Variables:
%%% mode: LaTeX
%%% TeX-master: t
%%% End:
